\section{Introduction}\label{sec:intro}

Filesystems intended for deployment in environments where persistent
storage is exposed to electromagnetic perturbation must, beyond the
crash-consistency guarantees of the FSCQ school~\cite{chen2015fscq},
provide a recovery contract that is empirically defensible. The
contract is empirically defensible only if (i) the perturbation
family is realizable by a documented physical model, (ii) the
recovery operator's behaviour under that family is stated
formally, and (iii) the formal statement is exposed to falsification
by a documented adversarial instrument. This report defends a
filesystem under all three conditions, after having retracted an
earlier statement that did not survive (iii).

\subsection{From radiation-only to electromagnetic resilience}

The v1 report of this lineage~\cite{desbrieres2026ftrfs} defended a
recovery calculus under a radiation-only model: the perturbation
family was the canonical single-event upset (SEU) and multi-bit
upset (MBU) regime of cosmic-ray and alpha-particle physics. The
soundness theorem (v1 Theorem IV.1) quantified over events whose
per-subblock symbol error count did not exceed the Reed--Solomon
correction capacity \(t = \lfloor (n-k)/2 \rfloor\).

On 2026-04-28, the companion fault injection operator
RadFI~\cite{desbrieres2026radfi} produced a counter-example: a
controlled perturbation pattern, formally an instance of the
operator's stated targeting algebra, that the v1 implementation
failed to recover and that did not return fail-closed. The
counter-example was not in violation of any hypothesis stated in
v1 Theorem IV.1 read literally; rather, it exhibited that the
hypothesis was insufficiently tight in two respects: the v1 model
assumed exclusively stochastic SEU, whereas RadFI implemented an
algebraically targeted perturbation, and the v1 implementation
lacked an explicit saturation boundary signalling.

The retraction is accepted in the Popperian
sense~\cite{popper1959logic}: a theorem retracted by empirical
falsification is the success state of the scientific method, not
its failure mode. The proper response is not to discard the
calculus but to enlarge the perturbation family covered by the
hypothesis, restate the recovery contract, and re-expose the
revised statement to falsification.

This report carries out that revision. The perturbation family is
restated as \emph{electromagnetic resilience} (EMR), comprising
both Family\,A stochastic perturbations (the v1 SEU/MBU regime,
preserved) and Family\,B adversarial perturbations (intentional
electromagnetic interference, high-power microwave, electromagnetic
pulse, treated under IEC\,61000-2-13 and
61000-4-36~\cite{iec61000-2-13,iec61000-4-36,radasky2004iemi}).
The recovery contract is split into two theorems: Theorem v2.1
covers the recovery of states perturbed by an event whose
per-subblock corruption is within the Reed--Solomon capacity;
Theorem v2.2 covers the saturation regime, formally requiring that
events exceeding the capacity yield fail-closed recovery with an
auditable journal entry flagging the uncorrectable
\((\text{block},\text{subblock})\) coordinates.

\subsection{Position with respect to neighbouring work}

BEAMFS v2 sits between three traditions. From the FTRFS
lineage~\cite{fuchs2015ftrfs} it inherits the principle of total
per-region Reed--Solomon protection and the operational scenario
of a CubeSat-grade filesystem; the present report's INLINE scheme
extends FTRFS by per-subblock granularity within each disk block,
yielding 128 symbol errors of correction capacity per 4\,KiB block
versus FTRFS's per-region granularity. From the formal
filesystem-verification tradition~\cite{chen2015fscq,klein2009sel4,sigurbjarnarson2016yggdrasil}
it inherits the discipline of stating a recovery contract as a
typed theorem with explicit hypotheses on the state space and the
operator; the present report does not claim machine-checked
verification, but it does state every hypothesis in a form that a
later mechanization can target. From the radiation hardness
assurance tradition~\cite{nasa-hdbk-4002a,esa-eccs60,milstd882e}
it inherits the convention of reporting recovery in terms of
correction capacity expressed in symbols and of bounding the
analysis by a saturation regime stated as a physical hypothesis.

\subsection{Contributions}

This report's contributions are:

\begin{enumerate}
\item An EMR threat model
  (\cref{sec:emr-threat}) decomposing electromagnetic perturbations
  into a stochastic Family\,A and an adversarial Family\,B, with
  explicit citation of the canonical physical and standardization
  literature for each family, and a transverse saturation boundary
  defined in Reed--Solomon symbol-count terms.
\item A revised filesystem state space and consistency relation
  (\cref{sec:state}), preserving the v1 invariants
  \textbf{I1}--\textbf{I5} on the layout and the recovery
  operator. The v1 mathematical scaffolding is preserved; only the
  perturbation family hypothesis is restated.
\item Two replacement theorems
  (\cref{sec:soundness}). \textbf{Theorem v2.1} states recovery
  under bounded EMR perturbation, with the
  termination/safety/coverage proof structure of v1 Theorem IV.1
  retained. \textbf{Theorem v2.2} states saturation observability,
  formally requiring that the recovery operator on a saturating
  event emits an auditable uncorrectable-flagged journal entry.
\item A three-tier descent from mathematics to algorithm to kernel
  implementation (\cref{sec:m2c}), with a traceability table
  pairing each operator step with the corresponding kernel
  function in the BEAMFS implementation, including the explicit
  INLINE read path, the autonomic in-place repair via
  \texttt{mark\_buffer\_dirty} and \texttt{sync\_dirty\_buffer},
  and the journal entry emission on uncorrectable saturation.
\item Empirical corroboration (\cref{sec:empirical}) including a
  byte-level disk corruption recovered byte-perfect by the INLINE
  read path (Palier 3), and three benchmark suites (Palier 4)
  reporting read cold latency, write+fsync latency, and INLINE
  read on the conformance fixture canary block, all reproducible
  from a tagged GPG-signed commit on a Yocto-built aarch64
  research image.
\end{enumerate}

\subsection{What this report does not claim}

This report does not claim machine-checked verification of the
recovery operator: the proofs in
\cref{sec:soundness,app:proof} are conventional pen-and-paper
arguments. It does not claim a write path for the INLINE scheme:
INLINE write is currently \texttt{-EOPNOTSUPP} and is part of the
v2.x roadmap (\cref{sec:limitations}). It does not claim
exhaustive coverage of all conceivable EMR events: Family\,B is
parametrized by IEC standards but real-world adversarial waveform
synthesis exceeds the present model's scope. Finally, this report
explicitly does not republish the v1 soundness theorem; the
present Theorems v2.1 and v2.2 supersede it.
